\documentclass[a4paper, 10pt]{article}
%\usepackage{fontspec}
%\setmainfont{Lato}
\usepackage{pgf}
\usepackage{eurosym}
\usepackage{graphicx}
\usepackage{wasysym}
\usepackage{hyperref}
\usepackage{listings}
\usepackage{pxfonts}
\usepackage{verbatim}
\usepackage{color}
\usepackage{xcolor}
\usepackage{wrapfig}
\usepackage{enumitem}
\usepackage{booktabs}
\usepackage{gensymb}
\usepackage{tabularx}
\usepackage{currfile}

\hypersetup{
    bookmarks=true,         % show bookmarks bar?
    unicode=true,          % non-Latin characters in Acrobat’s bookmarks
    pdftoolbar=true,        % show Acrobat’s toolbar?
    pdfmenubar=true,        % show Acrobat’s menu?
    pdffitwindow=true,     % window fit to page when opened
    pdftitle={Assessments},    % title
    pdfauthor={Paul Vesey},     % author
    pdfsubject={Advanced Graphics Assignment },   % subject of the document
    pdfcreator={},   % creator of the document
    pdfproducer={xelatex}, % producer of the document
    pdfkeywords={'Graphics' }, % list of keywords
    pdfnewwindow=true,      % links in new PDF window
    colorlinks=true,       % false: boxed links; true: colored links
    linkcolor=violet,          % color of internal links (change box color with linkbordercolor)
    citecolor=magenta,        % color of links to bibliography
    filecolor=red,      % color of file links
    urlcolor=blue           % color of external links
}

\setlength\parindent{0pt}
\begin{document}

\lstset{language=HTML,
				basicstyle=\small,
				breaklines=true,
        numbers=left,
        numberstyle=\tiny,
        showstringspaces=false,
        aboveskip=-20pt,
        frame=leftline
        }
				
\begin{table}%
	\begin{minipage}{0.4\textwidth}%
			\includegraphics[width=1\textwidth]{./img/LITlogo.jpg}
	\end{minipage}
	\qquad
	\centering
	\parbox{0.4\textwidth}{
		\begin{large}			
			\begin{tabular}{| r | l |} \hline
				Subject: & \textbf{Advanced Graphics}\\
								 & \textbf{\& Visualisation}\\
				Course: & \textbf{Interior Design Y3}\\
				Session: & \textbf{Autumn 2020}\\
				Lecturer: & \textbf{Paul Vesey \footnotesize{BEng, MIE, HDip}}\\
				Filename: & \footnotesize{\currfilename}\\
				\hline
			\end{tabular}
		\end{large}			
	}
\end{table}
\vspace{0.25cm}	
	
%     _            _                                  _     _____ 
%    / \   ___ ___(_) __ _ _ __  _ __ ___   ___ _ __ | |_  |___ / 
%   / _ \ / __/ __| |/ _` | '_ \| '_ ` _ \ / _ \ '_ \| __|   |_ \ 
%  / ___ \\__ \__ \ | (_| | | | | | | | | |  __/ | | | |_   ___) |
% /_/   \_\___/___/_|\__, |_| |_|_| |_| |_|\___|_| |_|\__| |____/ 
%                    |___/                                        

\newpage
\setcounter{page}{1}
\begin{center}
	\begin{figure}[ht]
		\centering
		\includegraphics[width = 6cm]{img/LITlogo.jpg}
		\label{fig:logoa3}
	\end{figure}
	\Large\textbf{Practical 3 - Intermediate Revit to Unity}\\
	\large\textbf{BA in Interior Design \& Technology}
\end{center}

In this assignment we are going to build on what has been previously accomplished. The key objective of this exercise to create a more realistic work-flow that will allow greater flexibility during the design process.  The key deliverables are:
\begin{enumerate}
	\item Separate the building structure from the furniture and fittings.
	\item Separate animation objects from the building structure.
	\item Animate a door in Unity.
\end{enumerate}

By separating the building elements we can then make edits to specific game assets quickly, without effecting other game assets.  In the case of furniture and fittings, this allows us to explore fit-out options quickly and with minimal disruption.\\

Separating doors and windows from the building structure will allow us to animate those objects within Unity.  Doors are normally animated to the open position when the user approaches during game play.  This is achieved by creating a trigger collider in the space adjacent to the door, which is then used to 'trigger' the door animation.\\

Before uploading your game to Moodle you should demonstrate your game to your lecturer.

\end{document}